\clearpage
\ihead{\thechapter. Festkörperphysik | \leftmark}
\chapter{Festkörperphysik}
\label{chap:festkoerperphysik}
\clearpage

\section{Einführung}

Die Festkörperphysik betrachtet große Systeme mit einer großen Zahl chemisch gebundener Atome. 

Diese befinden sich im festen Aggregatzustand, welcher zunächst genauer definiert und beschrieben wird.
Festkörper werden als Vielteilchen-Systeme mit $\sim\SI{e23}{Teilchen}$ beschrieben.
So können deren physikalischen Eigenschaften, Interaktionen, und zum Beispiel deren Deformation bestimmt werden.
Mit Methoden der Quantenmechanik und der Statistischen Physik werden die Näherungen und Modelle der Festkörperphysik erarbeitet.

Experimentell zeigt sich, dass die idealisierten Systeme in der Realität nicht exisitieren:
\begin{itemize}
	\item In jedem Gitter gibt es Abweichungen von einer perfekt periodischen Gitterstruktur.
	\item Jedes System weißt Defekte auf.
	\item In jedem System gibt es atomare Unreinheiten.
	\item In den Oberflächen der Festkörper gibt es oft starke Abweichungen von vereinfachenden Annahmen.
\end{itemize}
Die Forschung in der Festkörperphysik ist heutzutage eines der faszinierendsten und fortschrittlichsten Forschungsgebiete der Physik und bildet die Grundlage der 
\begin{itemize}
	\item Halbleitertechnologie (Mikro- und Nanoelektronik, Nanophotonik),
	\item Kommunikationstechnologie,
	\item und Quantentechnologie.
\end{itemize}

Die Charakterisierung eines Festkörpers findet über seine strukturellen, mechanischeln, thermischen, elektrischen, magnetischen und optischen Eigenschaften statt.

\paragraph{Vorläufige Anmerkungen}
Bei Normalbedingungen bzw. bei nicht-hochenergetischer Physik gibt es drei Aggregatzustände: fest, flüssig und gasförmig. 
Feste Stoffe lassen sich weiter aufteilen in kristalline und amorphe Festkörper:
\begin{itemize}
	\item kristalline Festkörper sind periodische Anordnungen kleiner 'Bauklötze'
	\item amorphe Festkörper besitzen eine feste, nicht periodische, dafür aber statistisch isotrope Anordnung. Eine bestimmte Nahordnung ist allerdings vorhanden. 
\end{itemize}
Kristalline Festkörper lassen sich weiter unterscheiden in
\begin{itemize}
	\item Polykristalline Festkörper, welche kleine kristalline Regionen beinhalten welche irregulär zueinander orientiert sind,
	\item Monokristalline Festkörper, in denen eine makroskopische Fernordnung herrscht. Die periodische Struktur bleibt über den gesamten Festkörper erhalten. Z.B. Edelsteine, künstliche Kristallzuchtbedingungen wie Silizium-Wafer oder III-IV-Halbleiter, und Metalle.
	\item Molekulare Kristalle, bei welchen die molekulare Struktur während der Kristallisation erhalten bleibt.
\end{itemize}

Zum Verständnis der Festkörperphysik ist es elementar zu wissen, welche Kräfte Atome zusammenhalten und welche strukturelle Ordnung dadurch realisiert wird.

\section{Kristalline Bindung}
\stepcounter{subsection}
Die elektrostatische Interaktion zwischen negativ geladenen Elektronen und positiv geladenen Atomkernen ist die fundamentale und einzige Ursache für den Zusammenhalt von Festkörpern. 
Die entstehenden Bindungen werden durch verschiedene Bezeichnungen definiert, welche für unterschiedliche Situationen verwendet werden. Diese sind wie folgt:
\begin{itemize}
	\item Van-der-Waals-Kräfte
	\item Kovalente Bindungen
	\item Metallische Bindungen
	\item Ionische Bindungen
	\item Wasserstoffbrücken
\end{itemize}
Die Ursache für die Unterschiede sind die Verteilung der Valenzelektronen und der Ionenrümpfen.

\subsubsection{Ionische Bindungen}
\paragraph{Ionisierungsenergie}
Die Ionisierungsenergie $I$ gibt an, wie viel Energie aufgebracht werden muss, um ein Elektron von einem Atom zu entfernen:
$$
\ce{Na} + I \to \ce{Na+} +  \ce{e-}
$$ 
Die Elektronenaffinität $A$ gibt an, wie viel Energie bei der Aufnahme eines Elektrons von einem Atom freigegeben wird:
$$
\ce{e-} + \ce{Cl} \to \ce{Cl-}+ A.
$$ 
Daraus ergibt sich Paulings Elektronegativitätsskala mit der atomspezifischen Elektronegativität
$$
K = \num{0,184}\cdot(I+A) \quad \text{in}\quad \si{\electronvolt}.
$$ 
Die Elektronegativitätsdifferenz zwischen zwei Elementen gibt Aufschluss über die ionische Charakteristik einer Bindung.
Je größer $I$ und $A$ eines Atoms sind, desto stärker ist die Tendenz während einer Bindung Elektronen anzuziehen.

\paragraph{Berechnung der Bindungsenergie}
Die Bindungsenergi $W_{ij}$ zwischen zwei einfach geladenen Ionen $i$ und $j$ ist
$$
W_{ij} = \pm \frac{\mathrm{e}^2}{4 \pi \varepsilon_0 r_{ij}}
$$ 
mit dem Abstand zwischen den beiden Ionen $r_{ij}$.
Das Vorzeichen hängt davon ab, ob es sich um gleichnamige oder ungleichnamige Ladungen handelt.
Echte Ionen bestehen aus räumlich ausgedehnten Kernen und Orbitalen. 
Eine Verkleinerung des Abstands zwischen zwei Ionen führt zu einer Überlappung dieser Orbitale.
Nach dem Pauli-Prinzip führt eine starke Überlappung der Orbitale zum Befüllen energetisch hoher Zustände.
Dadurch entsteht eine abstoßende Kraft zwischen den beiden Atomen.

\paragraph{Born-Mayer Abstoßungspotential}
Das Born-Mayer Abstoßungspotential lässt sich berechnen mit
$$
W_{\text{abs}} (r_{ij})= \frac{B}{r_{ij}^{n}}\quad n = 9\ldots 12,
$$ 
wobei sich die Parameter $B$ und $n$ entweder aus quantenmechanischen Berechnungen oder experimentellen Befunden, wie den Ergebnissen von Kompressibilitätstests, ergeben. 

Ein alternativer Ansatz für das Abstoßungspotential ist
$$
W_{\text{abs}} = \lambda e^{ - \frac{r_{ij}}{\rho}}.
$$ 
mit der Interaktionsstärke $ \lambda$ und der Interaktionsreichweite $\rho$. 

Ein Ion $i$, welches mit allen anderen Ionen $j$ des Körpers interagiert hat somit die Bindungsenergie
\stepcounter{equation}
\stepcounter{equation}
\begin{align}
	\label{eq:ii.2.3}
	W_{B,i} &= \sum_{j \neq i}^{} \left[ W_{ij} + W_{\text{abs}} (r_{ij}) \right]  \\
	\label{eq:ii.2.4}
	&= \sum_{j \neq i}^{} \left( \pm \frac{\mathrm{e}^2}{4 \pi \varepsilon_0 r_{ij}} + \frac{B}{r_{ij}^{n}} \right).
\end{align}

Mit dem Abstand zwischen zwei benachbarter Ionen $r$ lässt sich dies umschreiben zu
\begin{equation}
	\label{eq:ii.2.5}
	W_{B,i} = - \alpha \frac{\mathrm{e}^2}{4 \pi \varepsilon_0 r} + A_{n} \frac{B}{r^{n}}
\end{equation}
mit
$$
A_{n} = \sum_{j \neq i}^{} \frac{1}{P_{ij}^{n}} 
$$ 
und der \textbf{Madelung Konstante}
$$
\alpha = \sum_{j \neq i}^{} \left( \pm \frac{1}{P_{ij}} \right)
$$ 
und
$$
P_{ij} = \frac{r_{ij}}{r}.
$$ 
Das Vorzeichen in der Summe von $\alpha$ hängt analog davon ab, ob es sich um gleich- oder ungleichnamige Ladungen handelt. 
Die Bindungsenergie des gesamten Festkörpers mit $N$ Teilchen berechnet sich mit
$$
W_{B} = N W_{B,i}.
$$ 
und besitzt ein Minimum im Gleichgewichtszustand 
$$
\frac{\mathrm{d} W_{B}}{\mathrm{d} r}  \Big|_{r=r_{\mathrm{e}}} = N \frac{\mathrm{d}W_{B,i}}{\mathrm{d} r}  \Big|_{r= r_{\mathrm{e}}} = 0.
$$
Es folgt
\stepcounter{equation}
\begin{equation}
	\label{eq:ii.2.7}
	A_{n} B = \alpha \frac{\mathrm{e}^2}{4 \pi \varepsilon_0} \frac{r_{\mathrm{e}}^{n-1}}{r}.
\end{equation} 
Einsetzen von \ref{eq:ii.2.7} in \ref{eq:ii.2.5} und multiplizieren mit $N$ liefert die
\begin{important}
	\textbf{Bindungsenergie}
	\begin{equation}
		\label{eq:ii.2.8}
		W_{B} = - N \alpha \left[ 1 - \frac{1}{n} \right]  \frac{\mathrm{e}^2}{4 \pi \varepsilon_0 r_{\mathrm{e}}}.
	\end{equation}
\end{important}

Als Beispiel zur Berechnung von $\alpha$ betrachten wir das Natriumchlorid, dessen Werte für $P_{ij}$ in \autoref{tab:NaCl_P_werte} aufgelistet sind.
\begin{table}[H]
	\centering
	\caption{Anzahl der Nachbarn im Natriumchlorid-Kristall.}
	\begin{tabular}{cc} 
	\toprule
	 $P$ & Anzahl der Nachbarn \\ \midrule
	 $P_{i, i\pm 1}=1        $ & 6 \\ 
	 $P_{i, i\pm 2}=\sqrt{2}   $ & 12 \\ 
	 $P_{i, i\pm 3}=\sqrt{3}   $ & 8 \\ 
	 $P_{i, i\pm 4}=\sqrt{4}=2 $ & 6 \\ 
	 $\vdots $ &  $\vdots$ 	\\\bottomrule
	\end{tabular}
	\label{tab:NaCl_P_werte}
\end{table}
Es folgt
$$
\alpha = \frac{6}{\sqrt{1} } - \frac{12}{\sqrt{2} } + \frac{8}{\sqrt{3} } - \frac{6}{\sqrt{4} } + \ldots = \num{1,747}.
$$
Der Verlauf der Bindungsenergie im Natriumchlorid für die Atomabstände $r$ ist in \autoref{fig:vl14_abbildung_1} dargestellt.
\begin{figure}[H]
    \centering
    \def\svgwidth{.393\columnwidth}
    \subimport{./figures}{vl14_abbildung_1.pdf_tex}
	\caption{Verlauf der Bindungsenergie im Natriumchlorid für die Atomabstände $r$ mit Gleichgewichtsabstand $r_\mathrm e$.}
    \label{fig:vl14_abbildung_1}
\end{figure}

Weiter typische Werte für  $\alpha$ sind in Tabelle \ref{tab:typische_alpha} aufgelistet.
\begin{table}[H]
	\centering
	\caption{$\alpha$ für verschiedener ionischer Bindungen.}
	\label{tab:typische_alpha}
	\begin{tabular}{cc}
		\toprule
		Bindung & $\alpha$ \\\midrule
		\ce{NaCl} & \num{1,747} \\
		\ce{CsCl} & \num{1,762} \\
		\ce{ZnS}  & \num{1,638} \\
	\bottomrule
	\end{tabular}
\end{table}

Typische Bindungsenergien sind in Tabelle \ref{tab:Bindungsenergien} aufgelistet.

\begin{table}[H]
	\centering
	\caption{Bindungsenergien gängiger Ionenkristalle.}
	\label{tab:Bindungsenergien}
	\begin{tabular}{ccc}
		\toprule
		Bindung & \makecell{Bindungsenergie \\ pro Molekül} & \makecell{Bindungsenergie \\ pro \unit{\mol}} \\
		\midrule
		\ce{NaCl} & \SI{7,95}{\electronvolt} & \SI{764}{\kilo \joule}\\ 
		\ce{NaI} & \SI{7,10}{\electronvolt} & \SI{683}{\kilo \joule}\\ 
		\ce{KBr} & \SI{6,92}{\electronvolt} & \SI{663}{\kilo \joule}\\
		\bottomrule
	\end{tabular}
\end{table}

Solche Bindungen sind nicht-leitend und lösbar in polaren Lösemitteln. Die Bindung ist 	richtungsunabhängig.

\subsubsection{Kovalente Bindungen}
Rein kovalente Bindungen finden sich nur zwischen Atomen, die ähnliche bzw. gleiche Elektronegativität besitzen.

Eine Bindung wird typischerweise durch ein Elektronenpaar gebildet. 
Die Elektronen lokalisieren sich gemeinsam zwischen den beiden Atomen. 
Dadurch wird die Energie des Gesamtsystems minimiert, wobei die Spins {antiparallel} sind. 
Das Pauli-Prinzip modifiziert die Ladungsverteilung gemäß der Spinorientierung. 
Die Spin-abhängige Coulombinteraktion liefert die \textbf{Austauschinteraktion}, welche die kovalente Bindung erzeugt.

Bei kovalenten Bindungen herrscht, anders als bei ionischen Bindungen, eine starke Richtungsabhängigkeit.
Beispiele stark richtungsabhängiger Bindungen sind Diamanten, Silizium oder Germanium.

Die Bindungsenergie kovalenter Bindungen beträgt \SI{6}{\electronvolt} bis \SI{10}{\electronvolt} pro Atom.
Die Bindungen sind also sehr stark.

Stoffe mit kovalenten Bindungen sind Halbleiter, besitzen also nur ein moderate Leitfähigkeit.

Die Unterscheidung zwischen kovalenter und ionischer Bindung ist nicht binär, der Übergang ist wie die Elektronegativitätsdifferenz verschiedener Atome fließend. 
Rein ionische bzw. rein kovalente Bindungen sind somit Extrema, welche mit einer maximalen Differenz, bzw. gar keiner Differenz korrespondieren.
